%% kapa: exam=1 round=24 year=2013 period=2 form=A subject=accounting subject_ko=회계학 seq=30 number=30
%% source: Q-Net 감정평가사 기출문제 (cst00309) artlSeq=5073802 file=extracted/1008/5073802/2013년 제24회 감정평가사 1차시험 기출문제/2교시(A형).hwp
%% conversion: kordoc 4.24.51 -> tools/gen1.py (auto); review=proofread-sonnet; flags=none
\begin{question}{30}
(주)대한은 20×1년 1월 1일 (주)서울의 의결권주식 30\%(300주)를 주당 ￦1,500에 취득함으로써 유의적인 영향력을 행사할 수 있게 되어 관계기업투자주식으로 분류하였다. 취득 당시 (주)서울의 순자산 장부금액은 ￦900,000이었다. 취득 당시 (주)서울의 재고자산과 토지의 공정가치가 장부금액에 비해 각각 ￦100,000과 ￦200,000 더 높고 나머지 자산과 부채는 장부금액과 공정가치가 일치하였다. (주)서울의 재고자산은 20×1년에 모두 판매되었다. 20×1년도 (주)서울이 보고한 당기순이익은 ￦200,000이며 기타포괄이익은 ￦40,000이었다. (주)서울은 20×1년 12월 31일에 ￦30,000의 현금배당을 실시하였다. (주)대한이 지분법을 적용할 경우 20×1년도 말 관계기업투자주식은?
\begin{choices}
  \item ￦423,000
  \item ￦461,000
  \item ￦483,000
  \item ￦513,000
  \item ￦522,000
\end{choices}
\end{question}
\answer{3}
